% RESUMO - PT
\begin{resumo}
  \vspace{-15pt}
  
O monitoramento de áreas restritas por meio de imagens aéreas apresenta desafios importantes para a detecção automática de pessoas, em razão de fatores como reduzida escala dos alvos, variações de perspectiva, condições de iluminação e complexidade do ambiente observado. Nesse contexto, a visão computacional tem se mostrado uma alternativa promissora para apoiar processos de vigilância e identificação de presenças humanas em locais sensíveis. Este trabalho tem como objetivo avaliar modelos de detecção de pessoas aplicados ao monitoramento de áreas restritas em imagens aéreas, buscando identificar quais abordagens apresentam melhor desempenho nesse cenário. Para isso, serão analisados diferentes modelos de detecção de objetos baseados em aprendizado profundo, contemplando abordagens de uma etapa, duas etapas e arquiteturas mais recentes. A metodologia proposta consiste na aplicação e comparação desses modelos em uma base de imagens aéreas, com avaliação realizada por meio de métricas como precisão, recall, F1-score, Intersection over Union, mean Average Precision e tempo de inferência. Espera-se, com este estudo, contribuir para a análise comparativa de modelos de detecção de pessoas em cenários aéreos, fornecendo subsídios técnicos e acadêmicos para futuras aplicações no monitoramento de áreas restritas.


  Palavras-chave: Visão Computacional, Detecção de Pessoas, Imagens Aéreas, Monitoramento de Áreas Restritas, Detecção de Objetos 
\end{resumo}


% RESUMO - EN
\begin{resumo}[Abstract]
  \vspace{-15pt}
  
  \begin{otherlanguage*}{english}
  Monitoring restricted areas using aerial imagery presents significant challenges for the automatic detection of people due to factors such as small target scale, variations in perspective, lighting conditions, and the complexity of the observed environment. In this context, computer vision has shown promise as an alternative to support surveillance processes and the identification of human presence in sensitive locations. This work aims to evaluate people detection models applied to the monitoring of restricted areas in aerial imagery, seeking to identify which approaches perform best in this scenario. To this end, different object detection models based on deep learning will be analyzed, including one-stage, two-stage, and more recent architectures. The proposed methodology consists of applying and comparing these models to an aerial image database, with evaluation performed using metrics such as precision, recall, F1-score, Intersection over Union, mean Average Precision, and inference time. This study is expected to contribute to the comparative analysis of people detection models in aerial scenarios, providing technical and academic support for future applications in the monitoring of restricted areas.
  
  Keywords: Computer Vision, People Detection, Aerial Imagery, Restricted Area Monitoring, Object Detection
\end{otherlanguage*}
\end{resumo}
